\documentclass[11pt,letterpaper]{article}
\usepackage[top=.91in,bottom=.91in,left=1in,right=1in]{geometry}
\usepackage[T1]{fontenc}
\usepackage{lmodern,amsmath,amssymb,amsthm,microtype}
\usepackage[hidelinks]{hyperref}
\newtheorem{theorem}{Theorem}[section]
\newtheorem{lemma}[theorem]{Lemma}
\newtheorem{corollary}[theorem]{Corollary}
\newtheorem{proposition}[theorem]{Proposition}
\newcommand{\Z}{\mathbb Z}
\newcommand{\Q}{\mathbb Q}
\newcommand{\T}{\mathbb T}
\newcommand{\supp}{\operatorname{supp}}
\newcommand{\one}{\mathbf 1}
\newcommand{\abs}[1]{\left|#1\right|}
\newcommand{\norm}[1]{\left\|#1\right\|}
\setlength{\emergencystretch}{1.5em}
\allowdisplaybreaks[2]
\title{Recovering zero spectra and Dirichlet products\\from Li coefficients}
\author{}\date{November 27, 2026}
\begin{document}\maketitle
\begin{abstract}
Exponentiating a critical-line spectrum's complex Li coefficients at one positive time can erase integer multiplicities. We classify every invisible change and prove that each changes an ordinate of absolute value at most $\sqrt3/2$. The constant is sharp even for two singletons. At each exceptional time we determine the optimal central band, exactly which sparse probes suffice, the minimum number of altered atoms, and every equality case. For bounded complex weights on a fixed finite family of primitive Dirichlet characters we also establish uniform Lipschitz separation on every sufficiently long block of any fixed arithmetic progression after polynomial recalibration of bounded degree. For zeta and Dirichlet beta, 255 modes of stride 20 suffice to recover all integral multiplicities at time one with uniform log-magnitude error below $1/12$. The circle results hold for arbitrary admissible multisets. The arithmetic application uses the repository's critical-line input for Dirichlet zeros.
\end{abstract}

\section{What one time slice remembers}

For $\rho=\tfrac12+i\gamma$, the coordinate
\[
 z=1-\frac1\rho\in\T\setminus\{1\},\qquad
 \rho=\frac{1}{1-z},\qquad
 \gamma=\tfrac12\cot(\theta/2)\quad(z=e^{i\theta})
\]
identifies the critical line with the punctured circle. Multiplicity and the sign of $\gamma$ both matter. For a circle multiset $E$, write $m_E(z)$ for its nonnegative integer multiplicities. As in ProveIt \cite{repo}, call it \emph{admissible} when
\[
 \sum_zm_E(z)|1-z|^2<\infty,\qquad
 \lim_{\varepsilon\downarrow0}\sum_{|1-z|\ge\varepsilon}m_E(z)\Im z
 \quad\hbox{exists}.
\]
For $n\in\Z$ and $t>0$ define
\begin{equation}\label{eq:obs}
 \lambda_E(n)=\lim_{\varepsilon\downarrow0}
 \sum_{|1-z|\ge\varepsilon}m_E(z)(1-z^n),
 \qquad H_{E,t}(n)=\exp(-t\lambda_E(n)).
\end{equation}
The error in replacing $1-z^n$ by $-in\Im z$ is $O_n(|1-z|^2)$, which proves convergence. The time slice consists of all the modes $H_{E,t}(n)$, or equivalently of the probability measure with these Fourier coefficients obtained as the limit of the finite compound Poisson laws. Positive modes determine the negative ones by conjugation, and $H_{E,t}(0)=1$.

\begin{theorem}[Universal central band]\label{thm:central}
At any $t>0$, two admissible critical-line spectra with equal time slices and equal multiplicities at every signed ordinate $|\gamma|\le\sqrt3/2$ coincide at every height. The constant is sharp over all times. At $t=2\pi/\sqrt3$ the single zeros $\tfrac12+i\sqrt3/2$ and $\tfrac12-i\sqrt3/2$ have the same time slice.
\end{theorem}
The last pair corresponds to $e^{\pm i\pi/3}$, and its exponent differences lie in $i\sqrt3\,\Z$ for every mode. The upper bound, the optimal radius at each fixed $t$, and sparse alternatives to the interval follow in Section~\ref{sec:band}.

\paragraph{Relation to the repository.}
The merged article takes the upstream theorem for $\zeta$ and Dirichlet $L$-functions at the cited Lean revision as input. It derives admissibility of their Cayley multisets from the critical-line statement and convergence of the starred Li sums. The complex Li coefficients and their convergence are treated by Lagarias \cite{LagariasLi}. The repository's theorem \emph{Determining times} proves injectivity on admissible multisets exactly when $(\pi/t)^2\notin\Q$, and constructs a quadratic Gauss-sum collision at each remaining time. Its theorem \emph{Every time determines a primitive product} gives injectivity within the family of finite products of primitive Dirichlet $L$-functions using arithmetic rigidity. In the central-band theorem the competing spectrum may be any admissible multiset. We classify the full group of differences, determine all minimum collisions and characterize the local observations restoring injectivity. The proofs of these circle statements use no zero-location hypothesis.

\paragraph{Related work.}
Hasebe and Hotta study additive processes on the unit circle and discuss nonuniqueness of a marginal law's L\'evy measure \cite{hh}. Cyclotomic orbits also occur in Kadyan and Bhattacharjya's classification of integral mixed circulant graphs \cite{kb}, including the splitting of primitive shells into two character classes at conductor four. An additional condition here requires every Fourier coefficient to belong to the \emph{prescribed} lattice $2\pi i\Z/t$. Its spacing determines the least mass of an invisible change. The refinements here are relative to these results. Global bibliographic priority remains open.

\section{The integral kernel of the observation}

Write $\zeta_k=e^{2\pi i/k}$, $\mu_k=\{z:z^k=1\}$, $U_k=(\Z/k\Z)^\times$ and $\varphi(k)=|U_k|$. A collision consists of two distinct multisets with the same time slice. Set $b_z=m_E(z)-m_F(z)$. For a finite difference $b=\sum_z b_z\delta_z$, write $\widehat b(n)=\sum_z b_zz^n$. Its positive and negative parts are admissible multisets.

\begin{lemma}[Finite torsion reduction, after ProveIt]\label{lem:torsion}
Let $\beta=2\pi i/t$. Every collision difference is supported on some $\mu_q$, has mass zero, and satisfies $\widehat b(n)\in\beta\Z$ for every $n$. Moreover $\beta^2\in\Q$.
\end{lemma}
\begin{proof}
For $D=\lambda_E-\lambda_F$, equality of the time slices gives $D(n)\in\beta\Z$. Its central second difference is
\[
 u(n)=D(n+1)-2D(n)+D(n-1)=\sum_z b_z|1-z|^2z^n.
\]
The coefficients are absolutely summable. Approximate $u$ uniformly by a sum over finitely many $z$. Simultaneous recurrence of powers on a finite torus gives a positive $q$ for which those $z^q$ are all arbitrarily close to $1$: place sufficiently many successive powers in finitely many small boxes and subtract two indices landing in the same box. The uniform tail bound then gives $\sup_n|u(n+q)-u(n)|<|\beta|$. Since the values lie in $\beta\Z$, $u$ is $q$-periodic. Uniqueness of Fourier coefficients of finite measures, by density of trigonometric polynomials, gives $(z^q-1)|1-z|^2b_z=0$ for each $z$. This proves finite support.

Writing $s=\sum_z b_z$ gives $D(n)=s-\widehat b(n)$. Its real part vanishes, so Fourier inversion gives $b_z+b_{\bar z}=2s\one_{z=1}$. As $b_1=0$, $s=0$ and $D=-\widehat b$. Some $D(n_0)$ is a nonzero multiple of $\beta$, hence $\beta\in\Q(\zeta_q)$. The automorphism $\sigma_a:\zeta_q\mapsto\zeta_q^a$ for $a\in U_q$ sends $D(n_0)$ to $D(an_0)$, another nonzero integer multiple of $\beta$. Thus $\sigma_a(\beta)/\beta\in\Q^\times$ is a character of a finite group and takes values in $\{\pm1\}$. Every automorphism fixes $\beta^2$.
\end{proof}

At an exceptional time $\beta^2\in\Q_{<0}$, let $d<0$ be the fundamental discriminant of $\Q(\beta)$, set $f=|d|$, and let $\chi=(d/\cdot)$ be the primitive odd quadratic character of conductor $f$, extended by zero to nonunits. Its Gauss sum satisfies
\[
 G=\sum_{j\bmod f}\chi(j)\zeta_f^j,\qquad
 G^2=d,\qquad \sum_{j\bmod f}\chi(j)\zeta_f^{jn}=G\chi(n).
\]
These are the classical primitive Gauss sum identities, including the zero at nonunits \cite{davenport}. Since $\beta/G\in\Q^\times$, define coprime positive $A,B$ by $|\beta/G|=A/B$. In particular $|G|=\sqrt f$.

\begin{theorem}[All invisible changes]\label{thm:lattice}
Put $B_m=\sum_{j\bmod fm}\chi(j)\delta_{\zeta_{fm}^j}$ for $m\ge1$. At an exceptional time the complete group of differences having equal time slices is
\begin{equation}\label{eq:lattice}
 \mathcal K_t=\bigoplus_{m\ge1}\Z\,\frac{A}{\gcd(A,m)} B_m .
\end{equation}
For differences supported on $\mu_q$, retain precisely the generators with $fm\mid q$. The resulting rank is zero if $f\nmid q$, and the number of divisors of $q/f$ otherwise. At every other time the kernel is zero.
\end{theorem}
\begin{proof}
Take $Q=\operatorname{lcm}(q,f)$. For $a\in U_Q$ we have $\sigma_a(G)=\chi(a)G$. Hence $\widehat b(an)=\sigma_a(\widehat b(n))=\chi(a)\widehat b(n)$. Fourier inversion gives $b_{z^a}=\chi(a)b_z$. On the shell of primitive $k$th roots, $k\mid Q$, the power action of $U_Q$ is transitive. A nonzero coefficient there forces $\chi$ to be trivial on the kernel of $U_Q\to U_k$, equivalently $f\mid k$ by the definition of the conductor. The coefficients on such a shell have the form $v_k\chi(j)$ at $\zeta_k^j$ for $j\in U_k$, with $v_k\in\Z$.

Each $B_m$ is equivariant under this character. On its highest shell $k=fm$ the coefficient at $\zeta_{fm}$ is $1$, and its other shells have orders properly dividing $fm$. Successively removing shell coefficients in decreasing order gives a unique integral expansion $b=\sum_{fm\mid q}c_mB_m$. The argument applies to every integral array with this equivariance.

Write $j=u+f\ell$ and sum first over $\ell$ to obtain
\begin{equation}\label{eq:DFT}
 \widehat B_m(n)=
 \begin{cases}mG\chi(n/m),&m\mid n,\\0,&m\nmid n.\end{cases}
\end{equation}
Each $B_m$ has mass zero. Equality of time slices now requires $\widehat b(n)\in\beta\Z$. For $n\ge1$, the integer $\widehat b(n)/G=\sum_{m\mid n}mc_m\chi(n/m)$ lies in $(A/B)\Z$ exactly when $A$ divides it. At $n=m$ the new term is $mc_m\chi(1)=mc_m$. Induction gives $A\mid mc_m$ for all $m$. The same formula proves sufficiency, including negative modes by conjugation. Since $A\mid mc_m$ exactly when $A/\gcd(A,m)\mid c_m$, this proves the basis formula. Lemma~\ref{lem:torsion} handles the other times.
\end{proof}

The denominator $B$ drops out. Thus exceptional times with the same field and numerator $A$ have identical invisible changes. At $t=\pi$, $f=4$ and $A=1$, giving $B_1=\delta_i-\delta_{-i}$. At $t=\pi/3$ in the same field $A=3$, so the first two generators are $3B_1$ and $3B_2$, whereas $B_3$ itself is invisible.

\section{Low-ordinate probes and the sharp band}\label{sec:band}

For $k\ge3$ define the primitive ordinate set and its smallest absolute height by
\[
 \Gamma_k=\left\{\tfrac12\cot(\pi j/k):j\in U_k\right\},\qquad
 h_k=\min_{\gamma\in\Gamma_k}|\gamma|.
\]
These are potential zero locations available to a competing admissible spectrum.

\begin{theorem}[Exact supplementary data]\label{thm:probe}
At an exceptional time a time slice together with the multiplicities at all ordinates in $W\subseteq\mathbb R$ determines every admissible spectrum exactly when
\begin{equation}\label{eq:hit}
 W\cap\Gamma_k\ne\varnothing\qquad\text{for every positive multiple }k\text{ of } f.
\end{equation}
For $W=[-r,r]$ the exact threshold, including the endpoint, is $r_f=\tfrac12\tan(\pi/f)$. The positive ordinates $h_k$, one for each multiple of $f$, suffice. They accumulate only at zero.
\end{theorem}
\begin{proof}
A difference has shell coefficients $v_k\chi(j)$ as in Theorem~\ref{thm:lattice}. If it vanishes at one primitive $k$th root it vanishes on that entire shell. Conversely suppose $W$ misses the shell of order $k$. Assign $\chi(j)$ on that shell and zero elsewhere. This integral equivariant array is an integral combination of the $B_m$. Multiplication by $A$ makes it invisible by~\eqref{eq:lattice}. Its positive and negative parts give a collision agreeing on $W$.

A unit $j$ nearest to $k/2$ minimizes the absolute height. This yields
\begin{equation}\label{eq:heights}
 h_k=\begin{cases}
 \tfrac12\tan(\pi/(2k)),& k\text{ odd},\\
 \tfrac12\tan(\pi/k),&4\mid k,\\
 \tfrac12\tan(2\pi/k),&k\equiv2\pmod4,\quad k\ge6 .
 \end{cases}
\end{equation}
For odd $k$, choose $j=(k-1)/2$. If $4\mid k$, $k/2-1$ is odd and its common divisors with $k$ divide $2$. For $k\equiv2\pmod4$ the neighbors of $k/2$ are even, while $k/2-2$ is odd with common divisors dividing $4$. Each choice gives the positive height $h_k$.

If $f$ is even, $4\mid f$ by the fundamental discriminant classification, and the largest $h_k$ occurs at $k=f$. If $f$ is odd, the largest occurs at $k=2f$. Indeed the odd multiples have tangent argument at most $\pi/(2f)$, the multiples divisible by four have argument at most $\pi/(4f)$, and the remaining multiples have argument at most $\pi/f$. Both parity cases give the maximum $r_f$. Formula~\eqref{eq:heights} also gives $h_k\to0$. Finally $f\ge3$ and $r_f$ decreases with $f$. The pair in Theorem~\ref{thm:central} realizes $r_3=\sqrt3/2$, proving that theorem as well.
\end{proof}

\begin{corollary}[Ambiguity below the threshold]\label{cor:rank}
For $r>0$ the remaining group of invisible changes agreeing on $[-r,r]$ has finite rank
\[
 R_f(r)=\#\{k\in f\Z_{>0}:h_k>r\}
       =\frac{\pi}{2fr}+O(1)\qquad(r\downarrow0).
\]
\end{corollary}
\begin{proof}
Each missed shell allows one integral amplitude, and $A$ times each individual shell is in the kernel. Put $u=\arctan(2r)$. For even $f$, $k=fm$ is missed precisely when $m<\pi/(fu)$. For odd $f$, the tests on $m$ are $m<\pi/(2fu)$ for odd $m$, $m<\pi/(fu)$ for $4\mid m$, and $m<2\pi/(fu)$ for $m\equiv2\pmod4$. Their residue-class densities are $1/2$, $1/4$, $1/4$. Counting each class gives $\pi/(fu)+O(1)$ in either parity. Use $u=2r+O(r^3)$.
\end{proof}

Every admissible spectrum has finitely many atoms in a bounded interval of ordinates because $|1-z|$ is bounded away from zero there. The data in the universal band can therefore be supplied as a finite list of signed ordinates with multiplicities and the information that the list is complete.

\section{Minimum number of altered atoms}

The total variation $\|b\|_1=\sum_z|b_z|$ counts the atoms left on both sides after common multiplicities are cancelled. The positive and negative parts each contain $\|b\|_1/2$ atoms by Lemma~\ref{lem:torsion}.

\begin{lemma}[Norm of a character projection]\label{lem:norm}
Take $f\mid Q$ and $w\in\mu_Q$ of order $k$. The normalized average
\[
 P(w)=\frac{1}{\varphi(Q)}\sum_{a\in U_Q}\chi(a)w^a
\]
vanishes unless $f\mid k$. For $k=fh$ and $w=\zeta_k^v$, $v\in U_k$, it satisfies
\begin{equation}\label{eq:projection}
 P(w)=\frac{\chi(v)\mu(h)\chi(h)G}{\varphi(k)},\qquad
 |P(w)|\le\frac{\sqrt f}{\varphi(f)} ,
\end{equation}
where $\mu$ is the M\"obius function. Equality in the absolute bound occurs precisely at order $k=f$, or at $k=2f$ when $f$ is odd.
\end{lemma}
\begin{proof}
If $\chi$ is nontrivial on the kernel of $U_Q\to U_k$, the sum over that kernel vanishes. Otherwise $f\mid k$ and equal fibers reduce the average to $U_k$. Changing variables by $v$ extracts $\chi(v)$. Insert $\one_{(a,k)=1}=\sum_{e\mid(a,k)}\mu(e)$ in the numerator. Terms with $(e,f)>1$ vanish. For $e\mid k$ coprime to $f$, the inner sum is
\[
 \chi(e)\sum_{b\bmod k/e}\chi(b)e^{2\pi i b/(k/e)} .
\]
Then $e\mid h$. Summing over lifts modulo $f$ gives zero unless $k/e=f$, by a geometric series. The surviving term is $\mu(h)\chi(h)G$, interpreted as zero also when $(h,f)>1$. Since $\varphi(k)\ge\varphi(f)$, the bound follows. Equality requires $(h,f)=1$ and $\varphi(k)=\varphi(f)$, hence $\varphi(h)=1$. Precisely $h=1$ and, for odd $f$, $h=2$ qualify.
\end{proof}

\begin{theorem}[Exact distance]\label{thm:dist}
At an exceptional time the least total variation of a nonzero invisible difference is $A\varphi(f)$. Both $AB_1$ and $B_A$ attain it. The latter has no repeated atoms on either side. For packets of at most $M$ extra atoms over any fixed background, the observation is injective exactly when $M<A\varphi(f)/2$.
\end{theorem}
\begin{proof}
Choose $n$ with $\widehat b(n)\ne0$. Equation~\eqref{eq:DFT} and Theorem~\ref{thm:lattice} give $|\widehat b(n)|\ge A\sqrt f$. Since $\widehat b(an)=\chi(a)\widehat b(n)$, averaging over $a\in U_Q$ gives
\[
 |\widehat b(n)|=\abs{\sum_z b_zP(z^n)}
       \le \|b\|_1\,\sqrt f/\varphi(f).
\]
Each $B_m$ has $m\varphi(f)$ coefficients of modulus one, so the examples attain equality. Differences of bounded packets have norm at most $2M$. At the threshold take the positive and negative parts of a minimizing difference; the common background multiplies both time slices by the same nonzero sequence.
\end{proof}

\section{All collisions at the minimum distance}

Let $q$ be the least positive integer for which $\supp b\subseteq\mu_q$, and write $q=fm$. In particular the support generates $\mu_q$.

\begin{theorem}[Equality cases]\label{thm:eq}
If $f$ is even or $m$ is odd, the complete list at $q$ is
\[
            b=\pm(A/m)B_m,\qquad m\mid A.
\]
If $f$ is odd and $m$ even, set $s=\chi(2)$. Choose $X,Y\in\{-A,0,A\}$, not both zero, and retain the pairs for which
\[
 c=(Y-sX)/m\in\Z\setminus\{0\},\qquad e=2X/m\in\Z .
\]
The complete list at $q$ is then $b=cB_m+eB_{m/2}$. All minimizing differences are supported on $\mu_{fA}$ for even $f$ and on $\mu_{2fA}$ for odd $f$.
\end{theorem}
\begin{proof}
Suppose $\|b\|_1=A\varphi(f)$. Every nonzero $\widehat b(n)$ has modulus $A\sqrt f$ by the two bounds in the preceding proof. Equality there also forces $|P(z^n)|=\sqrt f/\varphi(f)$ at every $z\in\supp b$. Thus each $z^n$ has order $f$, or possibly $2f$ for odd $f$. The support generates $\mu_q$, so its image under the $n$th power map generates a group of order either $f$ or $2f$. That order is $q/\gcd(q,n)$. On the frequency side only the shell of order $f$ and, when $f$ is odd and $2f\mid q$, the shell of order $2f$ can occur.

On each frequency shell $\widehat b/G$ transforms by $\chi$ under multiplication of the index by units modulo $q$, and is thus determined by its value at a representative. The order-$f$ representative is $n=m$. By~\eqref{eq:DFT}, $\widehat B_m/G$ has value $m$ there and vanishes off that shell. When this is the only shell, inversion gives $b=\pm(A/m)B_m$, and integrality forces $m\mid A$.

In the second case put $X=\widehat b(m/2)/G$, $Y=\widehat b(m)/G$. The representatives of orders $2f$ and $f$ are $m/2$ and $m$. Formula~\eqref{eq:DFT} gives
\[
 \begin{array}{c|cc}
                 &m/2&m\\ \hline
 \widehat B_m/G  &0&m\\
 \widehat B_{m/2}/G &m/2&s\,m/2
 \end{array}
\]
and both transforms vanish off these shells. Solving at the two representatives and applying Fourier inversion yields the stated $c,e$. The integral triangular basis in Theorem~\ref{thm:lattice} gives their integrality; $c\ne0$ follows from minimality of $q$. Conversely the conditions give the Fourier lattice since $X,Y$ are multiples of $A$. On even exponents $j=2\ell$ at $\zeta_q^j$, the coefficient of $b$ equals $(sc+e)\chi(\ell)$. At odd exponents it equals $c\chi(j)$. Each parity accounts for $(m/2)\varphi(f)$ positions coprime to $f$. Hence
\[
 \|b\|_1=\tfrac12m\varphi(f)(|c|+|sc+e|)
 =\tfrac12\varphi(f)(|Y-sX|+|Y+sX|)
 =\varphi(f)\max(|X|,|Y|).
\]
The coefficient at $\zeta_q$ is $c\ne0$, so the least exponent is $q$. Finally $mc=Y-sX$ is a nonzero element of $\{\pm A,\pm2A\}$, implying $m\mid2A$.
\end{proof}

\paragraph{An odd-conductor example.}
At $t=2\pi/\sqrt3$, one has $f=3$, $A=1$, $s=-1$. At $q=3$ the minimum pair consists of the singleton multisets at $\zeta_3^{\pm1}$. At $q=6$ take $X=Y=1$; then $c=e=1$. The measure $B_2+B_1$ cancels on the primitive third roots and has coefficients $+1,-1$ at $\zeta_6,\bar\zeta_6$. This realizes the band endpoint. The equality classification gives no further minimum pairs at this time up to exchanging the sides and adding a common background.

\subsection*{Positivity and the number of possible spectra}
\begin{theorem}[Every time slice has a finite fiber]\label{thm:finitefiber}
At an exceptional time let $S_k$ be the total multiplicity at $+h_k$ and $-h_k$, for each $k$ divisible by $f$. These integers are determined by the sequence of moduli $|H_{E,t}(n)|$. Only finitely many $S_k$ are nonzero, and the number of admissible spectra giving this time slice satisfies
\[
 \#\{F:H_{F,t}=H_{E,t}\}\ \le\prod_{f\mid k}(1+S_k)<\infty.
\]
When $A=1$ equality can be realized for every prescribed finite list of positive $S_k$ with all other $S_k=0$. In particular, an admissible spectrum having a gap $|\gamma|\le r_f$ is determined by its time slice alone among all admissible competitors. The radius $r_f$ is sharp for this gap guarantee.
\end{theorem}
\begin{proof}
The modulus gives $\Re\lambda_E(n)=-t^{-1}\log |H_{E,t}(n)|$. Its central second differences are the Fourier coefficients of the finite measure
\[
 \tfrac12\sum_z |1-z|^2\bigl(m_E(z)+m_E(\bar z)\bigr)\delta_z.
\]
Fourier uniqueness recovers every total on a conjugate pair. The pairs at $\pm h_k$ are distinct for different $k$ since they belong to distinct primitive shells. They lie in $[-r_f,r_f]$, where admissibility permits only finitely many atoms. The multiplicity at $+h_k$ has at most $1+S_k$ choices. Once those choices are given, Theorem~\ref{thm:probe} determines the entire spectrum. When all $S_k=0$ this proves uniqueness without supplementary data.

For sharpness when $A=1$, on each prescribed shell put multiplicity $S_k$ at every point with $\chi(j)=+1$ and zero at the points with $\chi(j)=-1$. The two members of a conjugate pair have opposite character signs. The signed array $\chi(j)$ confined to each shell is invisible by Theorem~\ref{thm:lattice}. Independently transferring $0,1,\ldots,S_k$ units from the positive half to the negative half produces the stated product of distinct spectra. For any $A$, the pure-shell construction in the proof of Theorem~\ref{thm:probe} at a shell with $h_k=r_f$ gives two indistinguishable spectra both avoiding every smaller closed band.
\end{proof}

In fact, at every exceptional time every positive integer occurs as a fiber
size. Take $m=A^2$, $k=fm$. The array $V_k$ equal to $\chi(j)$ on the
primitive $k$th roots and zero elsewhere satisfies
\[
 V_k=\sum_{e\mid m,\ (e,f)=1}\mu(e)\chi(e) B_{m/e}.
\]
This is inclusion--exclusion on the exponents divisible by primes of $m$
outside $f$. For each squarefree $e\mid A^2$ one has $A\mid m/e$,
so every term belongs to the kernel. The single-shell construction in the
preceding proof now gives exactly $S_k+1$ spectra for any $S_k\ge0$.


For example, the first ordinates of both zeta and Dirichlet beta have
absolute value greater than $6$. Every finite product of these functions
therefore has a unique admissible spectrum at each positive time, even
when the competing spectrum is not assumed to come from an $L$-function.

\section{Finite blocks determine weighted primitive Dirichlet factors}
\label{sec:prog}

In this section assume the critical-line statement for primitive Dirichlet
$L$-functions in the pinned source. For a primitive character $\chi$,
including the character of conductor one, let $m_\chi(\rho)$ be the multiplicity
of a nontrivial zero. Write
\[
   z_\rho=1-1/\rho,\qquad
   \lambda_\chi(n)=\lim_{R\to\infty}
      \sum_{|\rho|\le R}m_\chi(\rho)(1-z_\rho^n).
\]
Here and below sums using $m_\chi$ run over distinct zeros. These are the
starred Li sums, with the orientation of the Cayley map fixed as displayed.
The classical zero count gives $\sum m_\chi(\rho)|1-z_\rho|^2<\infty$,
and the starred first sum converges. See Lagarias \cite{LagariasLi}.
The only possible limit point of the $z_\rho$ on the circle is $1$.
The repository proves identification from all positive modes for primitive
products at one time, and identification in natural density for individual
characters. Classical comparison theorems can use the zero divisors
directly. For instance, Banks proves a short-height-window discrepancy for
a pair of distinct primitive characters of the same conductor
(\href{https://arxiv.org/abs/2302.07073}{arXiv:2302.07073}, Thm.~1.1).
Our data consist of exponentiated Li coefficients at finitely many indices
in a prescribed progression.

Fix a nonempty finite set $\mathcal C$ of distinct primitive characters,
a stride $h\ge1$, a residue $0\le r<h$, and a nonnegative integer $d$.
For $b\in\mathbb C^{\mathcal C}$ put
\[
          A_b(k)=\sum_{\chi\in\mathcal C} b_\chi\lambda_\chi(r+hk).
\]
When the weights are positive, $\exp(-A_b)$ is a product of
heat coefficients at the indicated times. Complex weights allow comparison
and subtraction without choosing a branch of a complex power. A calibration
polynomial $P$ multiplies the observed sequence by $\exp(P(k))$.
It includes an unknown normalization and rotation.

\begin{theorem}[Uniform separation on a finite block]\label{thm:calibrated}
For each $M>0$ there exist an integer $W\ge1$ and a constant $c>0$,
depending only on $\mathcal C,h,r,d,M$, such that every interval $I$
of $W$ consecutive positive integers and every $b$ with $\|b\|_2\le M$ satisfy
\begin{equation}
 \inf_{\substack{P\in\mathbb C[X]\\ \deg P\le d}}
     \max_{k\in I}\left|\exp\big(-A_b(k)+P(k)\big)-1\right|
      \ \ge\ c\,\|b\|_2 .
 \label{eq:calibration}
\end{equation}
The same constants apply when a different calibration polynomial is fitted
on each interval.
\end{theorem}

For a catalogue of $m$ characters, every admissible $W$ satisfies
$W\ge m+d+1$. Indeed, on $W$ points the values of the polynomials of
degree at most $d$ form a subspace of dimension $\min(d+1,W)$.
If its codimension is smaller than $m$, the linear map $b\mapsto A_b|_I$
followed by the quotient has a nonzero kernel element. Scale it to have
norm at most $M$ and interpolate its values to make the left side zero.

\begin{lemma}[Rigidity of a finite weighted zero discrepancy]\label{lem:logr}
If $\sum_{\chi\in\mathcal C} b_\chi m_\chi(\rho)=0$ at all but finitely
many nontrivial zeros in the union of the zero sets, then $b=0$.
This lemma holds without a critical-line hypothesis.
\end{lemma}

\begin{proof}
Let $\xi_\chi$ be the entire order-one completion of $L(s,\chi)$, including
the factor $s(s-1)$ for the conductor-one character. Its zeros are exactly
the nontrivial zeros. The genus-one Hadamard formula gives
\[
  \frac{\xi_\chi'}{\xi_\chi}(s)
   =B_\chi+\sum_\rho m_\chi(\rho)
     \left(\frac{1}{s-\rho}+\frac{1}{\rho}\right).
\]
The paired summands converge normally away from their poles. Thus the
$b$-weighted sum of these logarithmic derivatives is a rational function,
a constant plus finitely many simple fractions. Write
$V(s)=\sum_\chi b_\chi L'(s,\chi)/L(s,\chi)$. The logarithmic derivative
of a gamma factor is $\tfrac12\psi((s+a_\chi)/2)$ with
$a_\chi\in\{0,1\}$. Its increment under $s\mapsto s+2$ is
$1/(s+a_\chi)$, by the gamma recurrence. The conductor contributes a
constant and the elementary factors contribute rational functions.
Consequently $V(s+2)-V(s)$ is rational. On the positive real axis it is
$O(2^{-s})$ by absolute convergence of the logarithmic Euler series.
A rational function with this decay vanishes identically.
Iterating the resulting period to the right gives $V=0$ on $\Re s>1$.

Uniqueness of Dirichlet series now yields
$\sum_\chi b_\chi\chi(p)=0$ at every prime $p$ coprime to the least
common multiple $q$ of the conductors. Every reduced class modulo $q$
contains a prime. Distinct primitive characters induce distinct characters
of $(\mathbb Z/q\mathbb Z)^\times$, so orthogonality gives $b_\chi=0$
for all $\chi$.
\end{proof}

\begin{proof}[Proof of Theorem~\ref{thm:calibrated}]
Set $j=\max(2,d+1)$ and $\Delta f(k)=f(k+1)-f(k)$.
The sequence $u_b=\Delta^j A_b$ has an absolutely convergent expansion
\begin{equation}
 u_b(k)=-\sum_z B_z z^r (z^h-1)^j(z^h)^k,\qquad
 B_{z_\rho}=\sum_{\chi\in\mathcal C}b_\chi m_\chi(\rho).
 \label{eq:push-prog}
\end{equation}
Indeed $|z^h-1|\le h|1-z|$ and $j\ge2$.
It follows that $u_b$ is uniformly almost periodic on $\mathbb Z$ and
depends continuously, in the supremum norm, on $b$.

Suppose $\exp(-u_b(k))=1$ for all $k\in\mathbb Z$.
The values of $u_b$ then lie in $2\pi i\mathbb Z$.
An almost period with error less than $2\pi$ is an exact period. Choose
such a positive period $v$. By uniqueness of Fourier coefficients of
finite measures on the circle, the measure in \eqref{eq:push-prog} after
grouping at $w=z^h$ is supported on $w^v=1$.
Choose an open arc about $1$ on which $z\mapsto z^h$ is injective.
There are finitely many atoms outside the arc. Remove their images and
the $v$th roots of unity. Above every remaining $w$, any contributing
atom is the sole preimage in that arc. Its coefficient in
\eqref{eq:push-prog} must vanish, and $w\ne1$, so $B_z=0$.
The discarded finite set has finite fibers. Thus $B_z$ is zero except
at finitely many $z$. Lemma~\ref{lem:logr} implies $b=0$.
In particular $u_b$ itself is nonzero whenever $b\ne0$.

For $\|a\|_2=1$ and $0\le s\le M$ consider the bounded sequences
\[
 g_{a,s}=
 \begin{cases}
  (\exp(-s u_a)-1)/s,& s>0,\\
  -u_a,&s=0.
 \end{cases}
\]
They are uniformly almost periodic. The map $(a,s)\mapsto g_{a,s}$ is
continuous into the space of bounded sequences with its supremum norm,
including at $s=0$ by the exponential series. Each sequence is nonzero
by the preceding argument. For a nonzero uniformly almost-periodic
sequence choose an index where its modulus exceeds $2\eta>0$.
The translates of this index by its $\eta$-almost periods have bounded
gaps, so some excursion above $\eta$ occurs in every interval of a
fixed length. The same length works with threshold $\eta/2$ throughout a
supremum-norm neighborhood of the sequence. A finite subcover of the compact
parameter space gives $L$ and $\eta_0>0$ such that for every $b$ with
$0<\|b\|_2\le M$ there is a $k$ in each block of length $L$ with
\begin{equation}
          |\exp(-u_b(k))-1|\ge\eta_0\|b\|_2 .
 \label{eq:recurrentdiff}
\end{equation}

Write $R(k)=\exp(-A_b(k)+P(k))$. Since $\Delta^jP=0$,
\[
     \prod_{i=0}^{j} R(k+i)^{(-1)^{j-i}\binom{j}{i}}
       =\exp(-u_b(k)).
\]
This rational function of $j+1$ inputs equals $1$ at $(1,\ldots,1)$.
On the closed polydisc $|R_i-1|\le1/2$, its first derivatives are bounded,
so its difference from $1$ is at most $C_j\max_i|R_i-1|$.
Take $W=L+j$ and choose $k$ from its first $L$ positions using
\eqref{eq:recurrentdiff}. Either some $|R(k+i)-1|>1/2$, or the
last display gives $\max_i|R(k+i)-1|\ge (\eta_0/C_j)\|b\|_2$.
Taking $c=\min(1/(2M),\eta_0/C_j)$ proves \eqref{eq:calibration}.
\end{proof}

\paragraph{Reading the estimate.}
Suppose two weighted lists from the fixed catalogue give coefficient
sequences whose ratio, after calibration, differs from $1$ by at most
$\varepsilon$ on a block of length $W$. If their weight difference has
norm at most $M$, that difference has norm at most $\varepsilon/c$.
At zero error the weights agree exactly. At a common known time, this
recovers every factor multiplicity of a primitive product. If time is
unknown, the identifiable quantities are time times multiplicity:
$P$ at time $at$ and $P^a$ at time $t$ yield exactly the same data for every
positive integer $a$. A fixed roster and a weight bound allow a common
block length. The general proof gives existence of $W,c$. An explicit instance for two
real characters follows below.

For any two individual finite mixtures one can take the union of their
catalogues and choose $M$ after their weights are specified.
Consequently even a sequence of blocks whose union has natural density
zero determines a finite mixture, provided their lengths tend to infinity
and the ratio approaches $1$ in relative measure on them after
polynomial calibration of fixed degree. Indeed, for distinct mixtures
every full subblock of length $W$ contains an error of at least
$c\|b\|_2$, a positive lower fraction. For example, testing the modes
$r+hk$ with $2^{j^2}\le k<2^{j^2}+j$ suffices in this asymptotic sense.
This conclusion controls relative errors in the Fourier coefficients;
they can themselves be very small.

\subsection{An explicit 255-mode theorem for zeta and Dirichlet beta}
Let $\beta(s)=L(s,\chi_{-4})$ be the Dirichlet beta function. Its coefficients
and those of $\zeta$ are real, so their Li coefficients are real.
Write $A_{u,v}(k)=u\lambda_\zeta(r+20k)+v\lambda_\beta(r+20k)$ for
$0\le r<20$. The following constants apply to all starting indices.

\begin{theorem}\label{thm:explicit255}
For every interval $I$ of $255$ consecutive positive integers, all real
$u,v$, and every real affine function $P$, one has
\[
  \max_{k\in I}|A_{u,v}(k)-P(k)|
          \ \ge\ \tfrac16\max(|u|,|v|).
\]
Consequently the logarithms of the heat-coefficient magnitudes of
$\zeta^a\beta^b$ at time one, at these 255 indices in any residue class
modulo 20, determine the nonnegative integer multiplicities $a,b$ even
in the presence of unknown affine drift. Uniform errors strictly below
$1/12$ in the log magnitudes permit at most one such integer pair.
\end{theorem}

\begin{proof}
We give the numerical inputs and bounds explicitly. The first two positive
ordinates of $\zeta$ lie in $(14.13,14.14)$ and $(21,22)$.
The first four for $\beta$ lie in $(6.02,6.03)$, $(10.24,10.25)$,
$(12.98,13.00)$ and $(16.34,16.35)$. The certified data come from
\href{https://flintlib.org/doc/acb_dirichlet/zeta-zeros.html}{FLINT's
indexed zeta zero enclosures} and the
\href{https://www.lmfdb.org/L/1/2e2/4.3/r1/0/0}{LMFDB entry for $\chi_{-4}$}
with its linked zero list and reliability statement. Both spectra are
symmetric and have no central zero.

For $\gamma>0$ put $w(\gamma)=z_{1/2+i\gamma}^{20}
=\exp(40i\arctan(1/(2\gamma)))$.
Choose as targets $w_\zeta,w_\beta$ from the first positive zero of
each function. They satisfy, for either target $w_*$,
\begin{equation}
 |1-w_*|^2>5/3,\qquad |w-w_*|>3/25
       \quad\text{at every other folded zero in the two spectra}.
 \label{eq:boxes}
\end{equation}
Here is the full gap check. In increasing order of positive ordinate,
the real parts of the first four folded frequencies in the union lie in
\[
 (-.987,-.984),\quad (-.373,-.369),\quad
 (.030,.034),\quad (.155,.157).
\]
The fourth belongs to $\zeta$. All remaining positive ordinates give real
part greater than $.340$. Conjugate frequencies have the same real part.
The absolute imaginary parts at the two targets exceed $.16$ and $.98$,
respectively, separating each from its conjugate. For $\gamma\ge10$ the
angle in the exponential decreases in $(0,\pi)$, so its cosine increases.
On the first interval the angle lies between $\pi$ and $3\pi/2$.
Thus rational endpoint evaluations of
$((\gamma^2-1/4+i\gamma)/(\gamma^2+1/4))^{20}$ prove all these bounds.

We also need a bound on the mass of all the other frequencies. With
multiplicity in each sum, the two completed functional equations give
\begin{equation}
 \sum_{\rho\text{ of }\zeta}\frac1{|\rho|^2}
 +\sum_{\rho\text{ of }\beta}\frac1{|\rho|^2}
   =2-2\log\pi+2\frac{\beta'(1)}{\beta(1)}<\tfrac12 .
 \label{eq:bothmass}
\end{equation}
Indeed for either real completion $\xi$ the difference of its logarithmic
derivatives at $1$ and $0$ equals $2\xi'(1)/\xi(1)$ by symmetry.
The Hadamard expansion gives the same difference as
$\sum 1/(\rho(1-\rho))=\sum 1/|\rho|^2$ by the critical-line input.
For $\zeta$ the resulting expression is $2+\gamma_E-\log(4\pi)$, and for
$\beta$ it is $\log(4/\pi)-\gamma_E+2\beta'(1)/\beta(1)$, where
$\gamma_E$ is Euler's constant. To check the strict bound, use
$\beta(1)=\pi/4>0.78$, $\log\pi>1.14$, and $\beta'(1)<0.30$.
For the last inequality the alternating series for $\beta'(1)$ is bounded
above by the five terms from $\log(3)/3$ through $\log(11)/11$.
Termwise differentiation at $1$ follows from uniform convergence of the
alternating derivative series on a real neighborhood of $1$.
The script \texttt{check\_window.py} verifies these elementary
inequalities and the frequency enclosures using rational arithmetic.

Write $D(k)=A_{u,v}(k+1)-2A_{u,v}(k)+A_{u,v}(k-1)$ and
$B=\max(|u|,|v|)$. Its Fourier series has an atom for each zero at
$w=z^{20}$ with coefficient $b_\rho z^r|1-w|^2$, where $b_\rho$ is the
corresponding scalar $u$ or $v$, counted with the multiplicity of the zero.
By $|1-z^{20}|\le20|1-z|$ and \eqref{eq:bothmass}, the sum of the absolute
coefficients is less than $200B$ when $B>0$. Choose the target belonging to
a scalar of modulus $B$. It has coefficient of absolute value greater
than $5B/3$ by \eqref{eq:boxes}.

Let $K(X)=(64^{-1}\sum_{j=0}^{63}X^j)^4=\sum_{j=0}^{252}a_jX^j$.
Its coefficients are nonnegative and sum to one. In the average
$\sum_{j=0}^{252}a_j\overline w_*^{\,j}D(k_0+j)$ the chosen atom contributes
its full coefficient times a unit factor. Every other frequency is
multiplied in modulus by at most $(25/96)^4$, by the geometric sum and
\eqref{eq:boxes}. Since $200(25/96)^4<1$, the average has absolute
value greater than $2B/3$. Some $|D(k_0+j)|$ therefore exceeds $2B/3$.
An affine function has zero second difference, and the absolute coefficient
sum of a second difference is $4$. At least one of the $255$ indices
from $k_0-1$ to $k_0+253$ has $|A_{u,v}-P|>B/6$.
The case $B=0$ is immediate. For two integer pairs, their nonzero difference
has $B\ge1$. Two fits with errors strictly below $1/12$ would violate the
inequality by the triangle inequality.
\end{proof}


\paragraph{Further questions.}
For competitors on specified primitive shells with bounded total mass, how many Fourier modes suffice to distinguish different observations at a prescribed error level? Can the constants in the arithmetic separation theorem be computed uniformly from the maximum conductor, stride and allowed recalibration degree? Both questions ask how accurately a finite experiment can approximate the identifiability proved here.

\begingroup\small
\begin{thebibliography}{9}
\bibitem{repo}
V. Reshetnikov, \emph{ProveIt: Generalized Riemann Hypothesis}, merged article, especially the subsections on Li heat kernels and determining times. Repository snapshot \texttt{4ccd0fdeda0b}, November 2026.
\href{https://github.com/VladimirReshetnikov/ProveIt/tree/4ccd0fdeda0b1cc316d2dfdccfdbcfea9d4b77fb/Analysis/GeneralizedRiemannHypothesis}{Pinned project} and
\href{https://github.com/VladimirReshetnikov/ProveIt/blob/4ccd0fdeda0b1cc316d2dfdccfdbcfea9d4b77fb/Analysis/GeneralizedRiemannHypothesis/Article/sections/li.tex}{Li section}.
The named upstream revision is \texttt{72fead102f66}.
\bibitem{LagariasLi}
J. C. Lagarias, Li coefficients for automorphic $L$-functions,
\emph{Ann. Inst. Fourier} \textbf{57} (2007), 1689--1740.
\href{https://doi.org/10.5802/aif.2311}{doi:10.5802/aif.2311}.
\bibitem{hh}
T. Hasebe and I. Hotta, Additive processes on the unit circle and Loewner chains,
\emph{Int. Math. Res. Not.} 2022, 17797--17848.
\href{https://doi.org/10.1093/imrn/rnab157}{doi:10.1093/imrn/rnab157}.
\bibitem{kb}
M. Kadyan and B. Bhattacharjya, Integral mixed circulant graphs,
\emph{Discrete Math.} \textbf{346} (2023), 113142.
\href{https://doi.org/10.1016/j.disc.2022.113142}{doi:10.1016/j.disc.2022.113142}.
\bibitem{davenport}
H. Davenport, \emph{Multiplicative Number Theory}, 3rd ed.,
revised by H. L. Montgomery, Springer, 2000.
\href{https://link.springer.com/book/9780387950976}{Publisher edition}.
\end{thebibliography}
\endgroup
\end{document}
